\chapter{Demi-vie, décroissance et stabilité}\label{ch:rs:half-life-decay-and-stability}

McLean et Pontiff (2016) ont suivi 97 variables que des études publiées avaient présentées comme prédictives de la section transversale des rendements boursiers,
au-delà des échantillons des études et au-delà de leur publication. Les rendements des portefeuilles étaient inférieurs de 26~\% hors échantillon et de 58~\% après
publication~; ils attribuent l'écart, 32~\%, à l'effet des investisseurs qui achètent ce qu'ils ont lu. Chaque signal dont dispose un desk se trouve
sur un chemin de ce type, et les questions de ce chapitre sont de savoir à quelle vitesse l'information d'un prédicteur s'épuise à mesure que l'horizon s'allonge, à quel point
elle est stable d'une année à l'autre, et comment distinguer un prédicteur mort d'un prédicteur simplement malchanceux. L'horloge de la première question
se compte en jours~; celle de la deuxième, en années~; la troisième est un problème de puissance statistique, et sa réponse est souvent que les données ne permettent pas de trancher.

\section{Décroissance selon l'horizon}

\begin{definition}[Courbe de décroissance de l'IC, autocorrélation du signal, demi-vie du signal]\label{def:rs:half-life-decay-and-stability:decay}
La \emph{courbe de décroissance de l'IC}\index{courbe de décroissance de l'IC} d'un prédicteur est son \omterm{def:rs:anatomy-of-a-predictor:ic}{coefficient d'information} avec le rendement de la seule période
$t + h$, en fonction du retard $h$. L'\emph{autocorrélation du signal}\index{autocorrélation du signal} est la corrélation de rang
en coupe transversale du prédicteur avec lui-même $k$ périodes plus tôt, moyennée sur les dates. La \emph{demi-vie du signal}\index{demi-vie du signal} est
le retard auquel elle tombe à la moitié~; pour un processus autorégressif d'ordre~1 de coefficient $\rho$ par période, elle vaut $-\ln 2/\ln\rho$ (la
demi-vie du livre~4, chapitre~4).
\end{definition}

Les deux courbes répondent à des questions différentes. La \omterm{def:rs:half-life-decay-and-stability:decay}{courbe de décroissance de l'IC} indique combien de temps l'information dure, et donc la période de détention~; l'\omterm{def:rs:half-life-decay-and-stability:decay}{autocorrélation du
signal} indique à quelle vitesse le prédicteur lui-même change, et donc le taux de rotation. Un livre pondéré par les rangs, reconstruit à chaque période, échange une fraction
d'environ $1 - \rho$ de lui-même, où $\rho$ est l'\omterm{def:rs:half-life-decay-and-stability:decay}{autocorrélation du signal} au retard un. Sur \texttt{firm.synthmkt}
(\cref{fig:rs:half-life-decay-and-stability:curves} et le tableau), quatre prédicteurs présentent quatre formes~:

\begin{center}
\small
\setlength{\tabcolsep}{4pt}%
\begin{tabular}{@{}lccccc@{}}
\toprule
prédicteur & IC, $h = 1$ & demi-vie, courbe & autocorr. & \omterm{def:rs:half-life-decay-and-stability:decay}{demi-vie du signal} & rotation\\
\midrule
retournement, 1 jour & 0.037 & $<$ 1 jour & $-0.04$ & --- & 104~\%\\
surprise, 60 jours & 0.010 & 20 jours (ajustement) & 0.983 & 39 jours & 1.7~\%\\
momentum 12-1 & 0.008 & aucune dans 60 jours & 0.995 & 130 jours & 0.53~\%\\
rapport valeur comptable/cours & 0.004 & aucune dans 60 jours & 0.9995 & 1\,386 jours & 0.05~\%\\
\bottomrule
\end{tabular}
\end{center}

\begin{omfigure}
\begin{tikzpicture}
\begin{axis}[width=11cm,height=5.6cm,xlabel={retard $h$ du rendement journalier unique (jours)},ylabel={IC de rang moyen},
  tick label style={font=\small},label style={font=\small},xmin=0,xmax=61,ymin=-0.01,ymax=0.04,grid=major,grid style={gray!25},
  yticklabel style={/pgf/number format/fixed,/pgf/number format/precision=2},scaled y ticks=false,legend columns=2,
  legend style={font=\footnotesize,at={(0.5,-0.25)},anchor=north,draw=none,fill=none,
  /tikz/every even column/.append style={column sep=8pt}},table/col sep=comma]
\addplot[omThm,thick,mark=*,mark size=1.4pt] table[x=h,y=reversal]{figdata/research/13-half-life-decay-and-stability/decay.csv};
\addplot[omDef,very thick,mark=square*,mark size=1.4pt] table[x=h,y=surprise]{figdata/research/13-half-life-decay-and-stability/decay.csv};
\addplot[black!60,thick,dashed,mark=triangle*,mark size=1.4pt] table[x=h,y=momentum]{figdata/research/13-half-life-decay-and-stability/decay.csv};
\addplot[gray,thick,dotted,mark=o,mark size=1.4pt] table[x=h,y=bp]{figdata/research/13-half-life-decay-and-stability/decay.csv};
\legend{retournement d'un jour,surprise de résultats (détenue 60 jours),momentum 12-1,rapport valeur comptable/cours}
\end{axis}
\end{tikzpicture}

\omcaption{\omterm{def:rs:half-life-decay-and-stability:decay}{Courbes de décroissance de l'IC} sur \texttt{firm.synthmkt}~: l'IC de rang de chaque prédicteur connu à la clôture du jour $t$ avec le rendement du
seul jour $t + h$, moyenné sur huit années de dates. Le retournement ne vit qu'un jour~; la dérive programmée de 60 jours de la surprise tombe à zéro à 60 jours~;
le momentum et la valeur décroissent trop lentement pour s'éteindre dans la fenêtre.}\label{fig:rs:half-life-decay-and-stability:curves}
\end{omfigure}

Le tableau comporte deux mises en garde. La courbe de la surprise n'est pas exponentielle~: chaque événement dérive à taux constant pendant 60 jours, et un signal
observé à un âge quelconque a une dérive restante qui décroît linéairement jusqu'à zéro, de sorte que la courbe est une droite à coude~; un ajustement exponentiel
annonce une demi-vie de 20 jours pour une durée de vie programmée de 60. Ajuster une forme revient à la supposer. Et les ICs journaliers des prédicteurs lents, 0.004 à 0.008,
sont proches de leur bruit d'échantillonnage~; leur décroissance ne se mesure pas sur des journées isolées, seulement sur des rendements plus longs et sur de nombreuses années. La lecture pratique
se fait sur la colonne de rotation~: le retournement d'un jour échange tout son livre chaque jour et paie les coûts correspondants (chapitre~27), le rapport valeur comptable/cours
un vingtième de pour cent.

\section{Stabilité d'une année à l'autre}

\begin{definition}[Stabilité de l'IC]\label{def:rs:half-life-decay-and-stability:stability}
La \emph{stabilité de l'IC}\index{stabilité de l'IC} d'un prédicteur est le comportement de son \omterm{def:rs:anatomy-of-a-predictor:ic}{coefficient d'information} sur des
périodes sans recouvrement~: la dispersion des ICs annuels (ou mensuels) autour de leur moyenne, comparée à la dispersion que leur seule erreur d'échantillonnage produirait.
\end{definition}

L'IC mensuel d'un signal est un nombre bruité, et son bruit n'est pas seulement une erreur d'échantillonnage en coupe transversale~: lorsque le prédicteur se charge sur
un facteur commun, le rendement de ce facteur le mois considéré fait bouger ensemble le rendement de chaque action et, avec lui, l'IC. Le momentum sur le marché synthétique
(constant pendant dix ans) présente un IC mensuel moyen de 0.023 et un écart type de 0.247~; ses moyennes annuelles vont de $-0.11$ à
0.18, et un intervalle à 95~\% obtenu par bootstrap par blocs pour la moyenne va de $-0.035$ à 0.071, à cheval sur zéro. La \omterm{def:rs:fundamental-analyst-and-event-features:sue}{surprise de résultats} a une moyenne de 0.040
avec un écart type de 0.041, et ses moyennes annuelles restent entre 0.02 et 0.06. Deux prédicteurs constants, dont l'un paraît stable et
l'autre ressemble à une succession de régimes~: la différence tient à l'exposition aux facteurs, non à la réalité.

\section{Ruptures structurelles}

\begin{definition}[Rupture structurelle, test CUSUM]\label{def:rs:half-life-decay-and-stability:break}
Une \emph{rupture structurelle}\index{rupture structurelle} est un changement, à une date donnée, dans les paramètres du processus qui engendre une série (ici,
la moyenne de l'IC d'un prédicteur). Le \emph{test CUSUM}\index{test CUSUM} de Brown, Durbin et Evans (1975) cumule les résidus récursifs standardisés
(chaque observation moins la moyenne de celles qui la précèdent, remise à l'échelle) et rejette la constance lorsque la trajectoire sort d'un couple de droites
qui s'écartent linéairement avec le temps, $\pm a(\sqrt K + 2k/\sqrt K)$ avec $a = 0.948$ à 5~\% pour $K$ résidus.
\end{definition}

L'autre démarche consiste à essayer toutes les dates~: on ajuste une moyenne avant et une moyenne après, on calcule la statistique F de la coupure, et l'on retient la plus grande sur
le milieu de l'échantillon (Andrews, 1993, a donné la théorie de cette statistique sup-F pour un changement en un point inconnu~; Bai et Perron, 1998,
l'ont étendue à plusieurs ruptures). \texttt{firm.decay} calcule sa p-valeur par permutation, ce qui conserve la distribution propre de la série (queues
épaisses comprises) et suppose seulement que, en l'absence de rupture, l'ordre des mois est indifférent.

Une rupture programmée permet de comparer. Avec \texttt{pead\_break} au début de la septième année, la dérive post-annonce du marché synthétique est
supprimée entièrement~: la moyenne mensuelle de l'IC de la surprise passe de 0.032 avant à 0.005 après (\cref{fig:rs:half-life-decay-and-stability:break}).
La statistique sup-F, 12.5, culmine 53 mois après le début de la série (la rupture est au mois 48, et les annonces anciennes continuent de dériver jusqu'à 60 jours après
elle), avec une p-valeur par permutation de 0.013. Le CUSUM ne franchit pas sa frontière~: sa trajectoire atteint 0.69 de celle-ci. Avec une dérive réduite de 70~\% au lieu d'être supprimée, aucun des deux tests ne voit rien sur ce marché (p-valeur sup-F de 0.93).
Un marché n'est qu'un tirage~; sur dix marchés aux graines différentes, sup-F détecte la suppression dans les dix, la réduction de 70~\% dans sept, et un changement qui
n'a jamais eu lieu dans un seul~; le CUSUM les détecte respectivement dans cinq, trois et un cas. Et le momentum, que rien n'a modifié, échoue aux deux tests~: son CUSUM franchit
la frontière le dernier mois, sa p-valeur sup-F vaut 0.023, après quelques années très favorables. Un test à 5~\% rejette un prédicteur stable une fois sur vingt par construction,
plus souvent lorsque la variance de l'IC change dans le temps, et un desk qui surveille cinquante prédicteurs chaque mois verra chaque année des ruptures
qui n'existent pas.

\begin{omfigure}
\begin{tikzpicture}
\begin{axis}[width=11cm,height=5.4cm,xlabel={année du marché synthétique},ylabel={IC moyen sur 12 mois},tick label style={font=\small},
  label style={font=\small},xmin=2.8,xmax=10,ymin=-0.02,ymax=0.08,grid=major,grid style={gray!25},
  yticklabel style={/pgf/number format/fixed,/pgf/number format/precision=2},scaled y ticks=false,legend columns=3,
  legend style={font=\footnotesize,at={(0.5,-0.25)},anchor=north,draw=none,fill=none,
  /tikz/every even column/.append style={column sep=8pt}},
  table/col sep=comma]
\addplot[black!55,thick,dashed] table[x=year,y=intact]{figdata/research/13-half-life-decay-and-stability/rolling_ic.csv};
\addplot[omDef,very thick] table[x=year,y=broken]{figdata/research/13-half-life-decay-and-stability/rolling_ic.csv};
\draw[omThm,thick] (axis cs:6,-0.02) -- (axis cs:6,0.08) node[pos=0.92,left,font=\footnotesize]{rupture programmée};
\legend{dérive intacte,dérive supprimée à partir du jour 1\,512}
\end{axis}
\end{tikzpicture}

\omcaption{Moyenne mobile sur douze mois de l'IC mensuel de la \omterm{def:rs:fundamental-analyst-and-event-features:sue}{surprise de résultats} sur \texttt{firm.synthmkt}, avec la dérive poste-annonce intacte et
avec celle-ci supprimée à partir du début de la septième année (jour 1\,512). Le test sup-F détecte la rupture (p-value 0.013)~; le CUSUM ne la détecte pas~; sur ce marché, une réduction de 70~\% n'est détectée par aucun des deux.}\label{fig:rs:half-life-decay-and-stability:break}
\end{omfigure}

\section{Décroissance après publication}

\begin{definition}[Décroissance après publication]\label{def:rs:half-life-decay-and-stability:postpub}
La \emph{décroissance après publication}\index{décroissance après publication} d'un prédicteur est la baisse de son rendement, ou de son IC, après que l'étude qui
le documentait a été publiée, par rapport à la période que l'étude aurait pu examiner.
\end{definition}

Trois des facteurs les plus étudiés la montrent dans la bibliothèque de données de Kenneth French (le chapitre ne stocke que des statistiques dérivées).
On les prend chacun de juillet 1963 à la fin de l'année de parution de leur article (Banz, 1981, pour la taille~; Fama et French, 1992, pour la valeur~; Jegadeesh et
Titman, 1993, pour le momentum), puis de là à juillet 2026~:

\begin{center}
\small
\begin{tabular}{@{}lccccc@{}}
\toprule
facteur (article) & moyenne mensuelle, avant & $t$ & moyenne mensuelle, après & $t$ & baisse\\
\midrule
SMB (1981) & 0.47~\% & 2.25 & 0.01~\% & 0.10 & 97~\%\\
HML (1992) & 0.42~\% & 3.12 & 0.19~\% & 1.16 & 55~\%\\
Mom (1993) & 0.85~\% & 4.76 & 0.37~\% & 1.54 & 56~\%\\
\bottomrule
\end{tabular}
\end{center}

Les baisses encadrent la moyenne de McLean et Pontiff, et les moyennes mobiles sur dix ans les montrent
(\cref{fig:rs:half-life-decay-and-stability:factors}). Pourtant, un test sup-F d'une rupture unique sur l'historique mensuel complet depuis 1963 ne rejette pour
aucun des trois à 5~\% (p-values de permutation 0.18, 0.15 et 0.053). Les rendements des facteurs sont assez bruités pour qu'une division par deux sur des décennies soit difficile
à dater statistiquement, et une date de publication n'est qu'une cause candidate parmi d'autres~: Chordia, Subrahmanyam et Tong (2014) ont constaté que les
rendements d'un portefeuille d'anomalies notoires ont à peu près été divisés par deux après la décimalisation, les actifs des hedge funds, l'encours vendeur à découvert
et la rotation expliquant la baisse. Une réplication ajoute une troisième lecture~: Hou, Xue et Zhang (2020) ont constaté que 65~\% des 452 anomalies ne pouvaient pas franchir une statistique $t$
de 1.96 une fois les microcapitalisations atténuées, et que même les anomalies répliquées étaient bien plus faibles que ce qui avait été initialement rapporté. Une partie de la
décroissance tient donc à la sélection de l'échantillon d'origine.

\begin{omfigure}
\begin{tikzpicture}
\begin{axis}[width=11cm,height=5.6cm,xlabel={année},ylabel={moyenne sur 120 mois (\% par mois)},tick label style={font=\small},
  label style={font=\small},xmin=1972,xmax=2027,ymin=-0.5,ymax=1.5,grid=major,grid style={gray!25},
  xticklabel style={/pgf/number format/1000 sep=},legend columns=3,
  legend style={font=\footnotesize,at={(0.5,-0.25)},anchor=north,draw=none,fill=none,
  /tikz/every even column/.append style={column sep=8pt}},table/col sep=comma]
\addplot[gray,thick] table[x=year,y=SMB]{figdata/research/13-half-life-decay-and-stability/factors.csv};
\addplot[omDef,thick] table[x=year,y=HML]{figdata/research/13-half-life-decay-and-stability/factors.csv};
\addplot[omThm,very thick] table[x=year,y=Mom]{figdata/research/13-half-life-decay-and-stability/factors.csv};
\draw[gray,dashed] (axis cs:1981,-0.5) -- (axis cs:1981,1.5);
\draw[omDef,dashed] (axis cs:1992,-0.5) -- (axis cs:1992,1.5);
\draw[omThm,dashed] (axis cs:1993,-0.5) -- (axis cs:1993,1.5);
\legend{SMB (publié en 1981),HML (1992),Mom (1993)}
\end{axis}
\end{tikzpicture}

\omcaption{Moyenne mobile sur 120 mois du rendement mensuel des facteurs américains de taille, de valeur et de momentum, à chaque mois de décembre à partir de 1973, avec les années des
articles qui les ont publiés (pointillés). Dérivée de la bibliothèque de données de Kenneth R. French (facteurs mensuels, fichier CRSP 202607).}\label{fig:rs:half-life-decay-and-stability:factors}
\end{omfigure}

\section{Mesurer la demi-vie d'une fiche de prédicteur}

Chaque \omterm{def:rs:anatomy-of-a-predictor:card}{fiche de prédicteur} (chapitre~6) comporte un champ d'horizon et de demi-vie, et les outils de ce chapitre le remplissent~: la \omterm{def:rs:half-life-decay-and-stability:decay}{courbe de décroissance de l'IC} à la
fréquence d'échantillonnage de la fiche, sa demi-vie ajustée avec un intervalle bootstrap (et l'avertissement de la surprise~: tracer la courbe avant de se fier à la
forme), l'\omterm{def:rs:half-life-decay-and-stability:decay}{autocorrélation du signal} et la rotation qu'elle implique, ainsi que la série mensuelle de l'IC avec ses \omterm{def:rs:half-life-decay-and-stability:break}{tests CUSUM} et sup-F, actualisée
chaque mois. La question la plus difficile est la dernière~: lorsque l'IC récent vaut la moitié de l'IC historique, le prédicteur est-il en train de mourir~? Deux périodes de $n$ mois
chacune, avec un IC mensuel de moyenne $\mu$ et d'écart type $\sigma$, donnent deux moyennes d'erreur type $\sigma/\sqrt n$~; la probabilité de
mesurer une division par deux alors que rien n'a changé, et celle de la mesurer alors que l'IC a réellement été divisé par deux, ne dépendent que de la statistique $t$ $\mu/(\sigma/\sqrt n)$ d'une période.

\begin{proposition}[La preuve d'une division par deux mesurée]\label{prop:rs:half-life-decay-and-stability:power}
Avec une probabilité a priori $\pi$ que l'IC ait réellement été divisé par deux, une division par deux mesurée (moyenne de la seconde période au plus égale à la moitié de la première) a pour probabilité a posteriori
$\pi p_1/(\pi p_1 + (1 - \pi)p_0)$ d'être une décroissance réelle, où $p_0$ et $p_1$ sont les probabilités de la mesure sans et avec décroissance~; pour des
moyennes normales, $p_1$ est proche de $\frac12$ et $p_0$ décroît avec la statistique $t$ de la période.
\end{proposition}

\begin{proof}
Règle de Bayes appliquée aux deux hypothèses. Avec décroissance, la seconde moyenne est centrée sur la moitié du centre de la première, avec le même bruit, et l'événement
«~seconde au plus égale à la moitié de la première~» est presque symétrique autour de ce centre, donc $p_1 \approx \frac12$~; sans décroissance, l'événement exige que le bruit
écarte les moyennes de la moitié de leur niveau, dont la probabilité décroît à mesure que le niveau croît par rapport au bruit.
\end{proof}

Avec $\pi = \frac12$ et des périodes de cinq ans (\cref{fig:rs:half-life-decay-and-stability:posterior}), un prédicteur affichant le bruit du momentum synthétique
(une statistique $t$ sur cinq ans de 0.72) qui mesure une division par deux a réellement décru avec une probabilité de 0.57, à peine plus que l'a priori~; un
prédicteur affichant le bruit du facteur momentum avant publication ($t = 1.93$), 0.72~; un prédicteur aussi stable que la surprise synthétique ($t = 7.6$), 0.999. Une baisse
constitue une preuve proportionnelle à la précision du prédicteur~; pour les prédicteurs bruités, la décision de l'abandonner doit reposer sur d'autres éléments (capacité, encombrement,
mécanisme) ou sur des échantillons plus longs.

\begin{omfigure}
\begin{tikzpicture}
\begin{axis}[width=10cm,height=5.4cm,xlabel={statistique $t$ sur cinq ans de l'IC},ylabel={P(décroissance réelle $\mid$ division par deux mesurée)},
  tick label style={font=\small},label style={font=\small},xmin=0,xmax=8.3,ymin=0.45,ymax=1.02,grid=major,grid style={gray!25},
  table/col sep=comma]
\addplot[omDef,very thick,mark=*,mark size=1.4pt] table[x=t,y=posterior]{figdata/research/13-half-life-decay-and-stability/posterior.csv};
\node[font=\footnotesize,anchor=west] at (axis cs:0.9,0.53) {momentum synthétique};
\node[font=\footnotesize,anchor=west] at (axis cs:2.1,0.70) {facteur momentum, 1963--1993};
\node[font=\footnotesize,anchor=east] at (axis cs:7.5,0.96) {surprise synthétique};
\end{axis}
\end{tikzpicture}

\omcaption{Probabilité a posteriori qu'un prédicteur dont la moyenne d'IC sur la seconde période de cinq ans est au plus égale à la moitié de la première ait réellement été divisé par deux, avec une probabilité a priori
d'un demi, en fonction de la statistique $t$ sur cinq ans de son IC (\cref{prop:rs:half-life-decay-and-stability:power}~; 200\,000 paires simulées par point).}\label{fig:rs:half-life-decay-and-stability:posterior}
\end{omfigure}

\section{Tutoriel~: dater la baisse d'un prédicteur}\label{tut:rs:half-life-decay-and-stability:tut}

\begin{tutorial}
\textbf{Objectif.} Mesurer les courbes de décroissance et les demi-vies de quatre prédicteurs synthétiques, tester une rupture programmée avec les \omterm{def:rs:half-life-decay-and-stability:break}{tests CUSUM} et sup-F, mesurer
trois facteurs réels avant et après publication, et calculer la preuve qu'apporte une division par deux mesurée. \textbf{État final~:}
\cref{fig:rs:half-life-decay-and-stability:curves,fig:rs:half-life-decay-and-stability:break,fig:rs:half-life-decay-and-stability:factors,fig:rs:half-life-decay-and-stability:posterior}~;
les deux tableaux.
\begin{tutsteps}
\item \textbf{Ajuster la courbe} par moindres carrés sur une grille de demi-vies, ICs négatifs compris.
\omcode{code/firm/decay/firm_decay.py}{44}{56}{Demi-vie par moindres carrés d'une courbe de décroissance de l'IC.}\label{lst:rs:half-life-decay-and-stability:fit}
\item \textbf{Le chemin CUSUM et sa frontière.}
\omcode{code/firm/decay/firm_decay.py}{85}{96}{Le CUSUM de Brown, Durbin et Evans pour une moyenne constante.}\label{lst:rs:half-life-decay-and-stability:cusum}
\item \textbf{Exécuter} \texttt{decay\_table()}, \texttt{stability()} avec et sans la rupture programmée, \texttt{factors()}, \texttt{power()} et
\texttt{fig\_decay.py}~; \texttt{rs\_fetch\_factors.py} une fois pour les statistiques françaises.
\end{tutsteps}
\textbf{Pour aller plus loin.} Programmer une décroissance progressive au lieu d'une rupture (une dérive qui baisse linéairement sur trois ans) et voir quel test la détecte
en premier~; refaire le tableau de la \omterm{def:rs:half-life-decay-and-stability:postpub}{décroissance après publication} avec l'échantillon de chaque article au lieu de 1963.
\end{tutorial}

\section{Construction~: la boîte à outils de décroissance}\label{bld:rs:half-life-decay-and-stability:decay}

\begin{build}
\bfield{Rôle} Les champs d'horizon, de demi-vie et de stabilité de chaque \omterm{def:rs:anatomy-of-a-predictor:card}{fiche de prédicteur}, mesurés de la même façon, ainsi que le suivi mensuel des
prédicteurs en production.
\bfield{Interface} \texttt{ic\_decay(signal, ret, horizons, dates)}, \texttt{half\_life\_fit}, \texttt{rank\_autocorr}, \texttt{half\_life\_ar},
\texttt{turnover}, \texttt{block\_bootstrap}, \texttt{cusum}, \texttt{sup\_f}, \texttt{sup\_f\_critical}, \texttt{sup\_f\_pvalue},
\texttt{decline\_power(ic, sd, n, ratio)}~; \texttt{firm.synthmkt} gagne \texttt{pead\_break} et \texttt{pead\_after} (désactivés par défaut).
\bfield{Règles} Une courbe de décroissance utilise des rendements sur une seule période à chaque retard, jamais des rendements cumulés~; chaque demi-vie est donnée avec la forme qu'elle suppose~;
la p-valeur d'un test de rupture est accompagnée du nombre de prédicteurs surveillés.
\bfield{Tests d'acceptation} \texttt{code/firm/decay/tests/}~: une demi-vie de signal de cinq jours programmée, retrouvée par l'ajustement de la courbe et par l'autocorrélation~; un
signal à un jour~; forme du bootstrap~; CUSUM silencieux sur le bruit et franchissant sa frontière sur une rupture programmée~; sup-F localisant la rupture, sa valeur critique simulée,
p-valeurs par permutation sur une rupture, sur le bruit et sur un bruit à queue épaisse~; les probabilités de puissance.
\bfield{Pour aller plus loin} Bai--Perron pour plusieurs ruptures~; un modèle bayésien de détection de ruptures avec une loi a priori sur la décroissance~; courbes de décroissance par tranche de capacité.
\end{build}

\begin{omsources}
\item R.~D.~McLean et J.~Pontiff, «~Does academic research destroy stock return predictability~?~», \emph{Journal of Finance} 71(1), 2016.
\item T.~Chordia, A.~Subrahmanyam et Q.~Tong, «~Have capital market anomalies attenuated in the recent era of high liquidity and trading
activity~?~», \emph{Journal of Accounting and Economics} 58(1), 2014.
\item K.~Hou, C.~Xue et L.~Zhang, «~Replicating anomalies~», \emph{Review of Financial Studies} 33(5), 2020.
\item R.~L.~Brown, J.~Durbin et J.~M.~Evans, «~Techniques for testing the constancy of regression relationships over time~», \emph{Journal of the
Royal Statistical Society B} 37(2), 1975.
\item D.~W.~K.~Andrews, «~Tests for parameter instability and structural change with unknown change point~», \emph{Econometrica} 61(4), 1993~;
J.~Bai et P.~Perron, «~Estimating and testing linear models with multiple structural changes~», \emph{Econometrica} 66(1), 1998.
\item R.~W.~Banz, \emph{Journal of Financial Economics} 9(1), 1981~; E.~F.~Fama et K.~R.~French, \emph{Journal of Finance} 47(2), 1992~;
N.~Jegadeesh et S.~Titman, \emph{Journal of Finance} 48(1), 1993~; Kenneth R. French Data Library.
\end{omsources}

\section{Exercices}

\begin{exercise}[$\star$]\label{exo:rs:half-life-decay-and-stability:1}
Calculer la demi-vie d'un signal dont l'autocorrélation de rang quotidienne à un retard vaut 0.98, ainsi que le taux de rotation quotidien d'un portefeuille qui le détient.
\end{exercise}

\begin{exercise}[$\star$]\label{exo:rs:half-life-decay-and-stability:2}
Calculer le recul d'un facteur qui a rapporté 0.85~\% par mois avant publication et 0.37~\% après, et comparer ce recul
à la moyenne de McLean et Pontiff.
\end{exercise}

\begin{exercise}[$\star$]\label{exo:rs:half-life-decay-and-stability:3}
Avec $K = 94$ résidus récursifs, calculer la frontière du CUSUM au seuil de 5~\% au premier résidu et au dernier résidu.
\end{exercise}

\begin{exercise}[$\star\star$]\label{exo:rs:half-life-decay-and-stability:4}
Chaque événement dérive à un rythme constant pendant 60 jours après son annonce, et le signal est observé à un âge
aléatoire uniforme compris entre 0 et 59 jours. Montrer que la part des signaux qui dérivent encore au jour $t + h$ vaut $(60 - h)/60$, et expliquer pourquoi la \omterm{def:rs:half-life-decay-and-stability:decay}{courbe de décroissance de l'IC}
tombe ensuite linéairement à zéro à 60 jours, plutôt que de se réduire de moitié à taux fixe.
\end{exercise}

\begin{exercise}[$\star\star$]\label{exo:rs:half-life-decay-and-stability:5}
Un IC mensuel a pour moyenne 0.04 et pour écart type 0.10. Calculer la statistique $t$ sur cinq ans et estimer
la probabilité a posteriori que la décroissance soit réelle lorsque l'IC mesuré est divisé par deux, avec une probabilité a priori d'un demi.
\end{exercise}

\begin{exercise}[$\star\star$]\label{exo:rs:half-life-decay-and-stability:6}
Expliquer pourquoi un prédicteur qui charge sur un facteur commun a un IC mensuel plus bruité qu'un prédicteur qui n'y charge pas, même lorsque l'information réelle est la même.
\end{exercise}

\begin{exercise}[$\star\star\star$]\label{exo:rs:half-life-decay-and-stability:7}
\emph{Programmation.} Avec \texttt{rs\_decay.detection}, mesurer à quelle fréquence le sup-F et le CUSUM détectent une réduction de 50~\% de la dérive à l'année 7 sur les
dix graines, et la situer entre la suppression du chapitre et la réduction de 70~\%.
\end{exercise}

\begin{exercise}[$\star\star\star$]\label{exo:rs:half-life-decay-and-stability:8}
\emph{Trouver la faille.} «~L'IC de notre signal de momentum sur les 24 derniers mois représente le tiers de sa moyenne sur 10 ans, et le CUSUM a franchi sa frontière le
mois dernier~; nous le désactivons.~»
\end{exercise}

\section{Problème~: est-il mort, ou a-t-il eu de la malchance~?}

\begin{problem}[{Problème du week-end --- une baisse, pesée}]\label{pb:rs:half-life-decay-and-stability:1}
Les quatre prédicteurs synthétiques du chapitre, la surprise dont la dérive est supprimée à l'année 7, et les facteurs français.

\medskip\noindent\textbf{Partie I --- Horizons.}
\begin{enumerate}
\item Donner l'IC à un jour et l'\omterm{def:rs:half-life-decay-and-stability:decay}{autocorrélation du signal} de chaque prédicteur.
\item Expliquer pourquoi le portefeuille de retournement se renouvelle entièrement chaque jour.
\item Donner la demi-vie qu'un ajustement exponentiel attribue à la surprise, et rappeler sa demi-vie programmée.
\item Expliquer pourquoi la décroissance du momentum et celle de la valeur ne peuvent pas être mesurées sur des jours isolés.
\end{enumerate}

\medskip\noindent\textbf{Partie II --- Années.}
\begin{enumerate}[resume]
\item Donner la moyenne et l'écart type de l'IC mensuel pour le momentum et pour la surprise.
\item Donner l'intervalle de bootstrap de l'IC moyen du momentum.
\item Expliquer pourquoi l'IC du momentum est beaucoup plus bruité.
\item Dire ce que les \omterm{def:rs:half-life-decay-and-stability:break}{tests CUSUM} et sup-F indiquent au sujet du momentum, dont le simulateur maintient la valeur constante.
\end{enumerate}

\medskip\noindent\textbf{Partie III --- Ruptures.}
\begin{enumerate}[resume]
\item Donner l'IC moyen de la surprise avant et après la suppression programmée.
\item Dire où le sup-F place la rupture, avec quelle statistique et quelle p-valeur. Expliquer pourquoi elle ne se situe pas exactement au mois 48.
\item Dire ce que révèle le CUSUM.
\item Dire ce qui se passe avec une réduction de 70~\% à la place, sur ce marché et sur dix marchés.
\item Donner les baisses des trois facteurs après publication, ainsi que les p-valeurs du sup-F sur l'ensemble de leurs historiques.
\end{enumerate}

\medskip\noindent\textbf{Partie IV --- Le verdict.}
\begin{enumerate}[resume]
\item Donner $p_0$ et $p_1$ pour un prédicteur qui présente le bruit du momentum synthétique.
\item Énoncer le \emph{résultat nommé}~: la probabilité qu'un signal dont l'IC mesuré a diminué de moitié sur cinq ans ait réellement décru, pour les trois
précisions du chapitre.
\item Donner la statistique $t$ qui rend une division par deux mesurée convaincante à 90~\%.
\item Dire quels autres éléments il faudrait rechercher avant de désactiver un prédicteur bruité.
\item Expliquer comment fixer une règle de suivi pour cinquante prédicteurs afin que les fausses alarmes restent rares.
\item Dire quelles autres explications que les transactions des investisseurs peuvent rendre compte d'une \omterm{def:rs:half-life-decay-and-stability:postpub}{décroissance après publication}.
\item Dire en une phrase quand une décroissance constitue une preuve.
\end{enumerate}
\end{problem}

\section{Questions d'entretien}

\begin{interviewq}[$\star$ \iqroles{chercheur}]\label{iq:rs:half-life-decay-and-stability:1}
Quelle est la demi-vie d'un signal, et quel lien entretient-elle avec le taux de rotation~?
\end{interviewq}

\begin{interviewq}[$\star\star$ \iqroles{chercheur}]\label{iq:rs:half-life-decay-and-stability:2}
Comment mesureriez-vous la vitesse à laquelle l'information d'un prédicteur décroît avec l'horizon~?
\end{interviewq}

\begin{interviewq}[$\star\star$ \iqroles{chercheur, trader}]\label{iq:rs:half-life-decay-and-stability:3}
Votre meilleur signal a connu une mauvaise année. Comment décidez-vous s'il est défaillant~?
\end{interviewq}

\begin{interviewq}[$\star\star$ \iqroles{chercheur}]\label{iq:rs:half-life-decay-and-stability:4}
Pourquoi les anomalies s'affaiblissent-elles après publication~? Quelles sont trois explications possibles, et comment les distinguer~?
\end{interviewq}

\begin{interviewq}[$\star\star$ \iqroles{chercheur}]\label{iq:rs:half-life-decay-and-stability:5}
Comment comparez-vous les \omterm{def:rs:half-life-decay-and-stability:break}{tests CUSUM} et sup-F pour détecter un changement de la moyenne d'un signal~?
\end{interviewq}

\begin{interviewq}[$\star\star\star$ \iqroles{chercheur}]\label{iq:rs:half-life-decay-and-stability:6}
Comment obtiendriez-vous la probabilité que la moyenne de l'IC du signal sur sa seconde période de cinq ans soit au plus la moitié de celle de la première, lorsque rien n'a changé, en fonction de la
statistique $t$ de l'IC~?
\end{interviewq}
